\chapter{Methodology}

\section{Baseline System and Proposed Extensions}


{\small

Hydra builds the 3D map and publishes local graph context in real time. The VLM pipeline asynchronously consumes this context together with RGB-D frames, produces object/risk observations, and publishes them back. A Hydra-side fusion module validates and integrates these observations into the global Dynamic Scene Graph.
}

\section{Hardware Setup and Input Data Preparation}

This section describes the hardware setup used for data collection and the preparation of input data for the proposed risk-annotated scene graph pipeline. In a complete deployment scenario, the system is expected to operate on a mobile robot equipped with an RGB-D camera and sufficient onboard computing resources to run the full perception, mapping, and risk-assessment pipeline in real time. Such a setup would allow sensor data acquisition, semantic interpretation, scene graph construction, and risk annotation to be performed directly on the robot.

For the current development and evaluation phase, however, the available hardware is distributed across three separate components: the RGB-D sensing unit, the mobile robotic platform, and the external processing unit. The RGB-D sensor used in this work is an Intel RealSense D455 camera, which provides synchronized color and depth information. The mobile platform is a Unitree Go1 quadruped robot, which provides the physical mobility required for collecting data in indoor environments. The processing unit is an NVIDIA Jetson Thor, which is intended to execute the computationally demanding components of the pipeline, such as visual-language-model-based reasoning and semantic scene graph processing.

In the present setup, the Unitree Go1 does not provide the required hardware interface to power and operate the NVIDIA Jetson Thor directly as an onboard edge-computing device. Therefore, the system is used primarily for data collection, while the computationally demanding components of the pipeline are executed later on the NVIDIA Jetson Thor. The Intel RealSense D455 is mounted on the Unitree Go1 using a 3D-printed frame and connected to the robot’s onboard computer through USB. During operation, the RGB-D sensor data, robot odometry, and coordinate transform information are transmitted to a local computer through Wi-Fi or Ethernet and recorded as a ROS bag. The recorded data can then be replayed at real-time speed for offline processing, testing, and debugging. This workflow is useful during development because the perception and risk-assessment pipeline can be evaluated repeatedly without physically operating the robot for every debugging cycle.




\section{Pre-Processing}

